% Источники: exam/materials/моделирование.pdf, разделы 47, 15 и 17; GitHub HumsterProgrammer/iu9-notes/8sem/Моделирование/Моделирование_лекция-04_19-02-26.md; GitHub HumsterProgrammer/iu9-notes/8sem/Моделирование/Моделирование_лекция-07_05-03-26.md; GitHub HumsterProgrammer/iu9-notes/8sem/Моделирование/Моделирование_лаб-3_02-04-26.md.
\section{Методика построения имитационной стохастической модели.}

В имитационном моделировании участвуют несколько специалистов: специалисты предметной области, математики и IT-специалисты.

\subsection*{Методика построения}

Используется общая методика построения имитационной модели: постановка задачи, структурная схема, концептуальная модель, формализация, программная реализация, отладка, серия прогонов, проверка качества и модификация. Отличие — на этапе концептуального проектирования модель классифицируется как \textbf{стохастическая}, поэтому случайные входы, параметры и выходы описываются распределениями.

\subsection*{Концептуальная модель: стохастическая модель}

\textbf{Стохастическая модель} — такая модель, входными и выходными данными которой являются случайные величины, а преобразование входных данных в выходные осуществляется посредством случайной функции (процесса).

Если стохастическая — то \textbf{статистическая}, \textbf{вероятностная} или \textbf{феноменологическая}:

\begin{itemize}
  \item \textbf{Вероятностная модель} — основана на основной задаче теории вероятностей (поиск вероятности в схеме Бернулли).
  \item \textbf{Статистическая модель} — основана на одной из трёх задач математической статистики:
    \begin{itemize}
      \item точечное оценивание параметров распределения характеристики объекта;
      \item интервальное оценивание (построение доверительных интервалов);
      \item проверка статистических гипотез.
    \end{itemize}
  \item \textbf{Феноменологическая модель} — построена по результатам конкретного (уже оконченного) эксперимента.
\end{itemize}

\textbf{Статистическая гипотеза} $H$ — любое утверждение относительно функции распределения $F(x)$ некоторой СВ $x$: о типе ФР, значениях её параметров и т.д.
\textbf{Нулевая гипотеза} $H_0$ — гипотеза, справедливость которой проверяется в ходе эксперимента.
\textbf{Альтернативные гипотезы} $H_1$ — отклонения от $H_0$.
\textbf{Статистический критерий} — совокупность правил, позволяющих по полученной выборке принять $H_0$ или отвергнуть гипотезу $H_1$.

Гипотеза \textbf{простая}, если её условиям удовлетворяет единственная ФР $F(x)$; иначе — \textbf{сложная}.

\subsection*{Что обязательно задаётся в стохастической имитационной модели}

\begin{itemize}
  \item случайные входные величины и процессы, их законы распределения или эмпирические распределения;
  \item генераторы псевдослучайных чисел и способы моделирования случайных величин;
  \item число прогонов, начальные условия, период разгона и правило остановки;
  \item оцениваемые характеристики: средние, дисперсии, вероятности событий, квантили;
  \item доверительные интервалы и критерии проверки адекватности.
\end{itemize}

Результат стохастической имитации — не одно число, а статистическая оценка характеристики модели с указанием точности. Поэтому после моделирования проводят серию независимых прогонов и применяют методы точечного, интервального оценивания и проверки гипотез.

\clearpage
